% Keep this file next to textbookchapter.cls.
% Build with: latexmk -pdf textbookchapter-example.tex
% On Overleaf, upload both source files and select this file as Main document.
% Run LaTeX at least twice when building by hand so the contents lists populate.
\documentclass[11pt,letterpaper,oneside,openany]{textbookchapter}
\usepackage{booktabs}

\hypersetup{pdftitle={Foundations of Linear Algebra},pdfauthor={Sample textbook}}

% To change the palette or the default labels, edit these lines and uncomment.
% \definecolor{ChapterAccent}{HTML}{007D9E}
% \definecolor{ChapterPale}{HTML}{E5F7FD}
% \definecolor{ChapterInk}{HTML}{171717}
% \renewcommand{\chaptercontentsname}{In This Chapter}
% \renewcommand{\chapterintroductionname}{Overview}
% \setcounter{tocdepth}{1} % Hide subsections in the book-wide contents only.

\begin{document}

\frontmatter
\hypersetup{pageanchor=false}
\begin{titlepage}
  \vspace*{.19\textheight}
  \noindent\textcolor{ChapterAccent}{\rule{34bp}{5bp}}\par
  \vspace{17bp}
  {\sffamily\raggedright\fontsize{36bp}{42bp}\selectfont
   \noindent Foundations of\\Linear Algebra\par}
  \vspace{18bp}
  {\sffamily\large\noindent Sample textbook\par}
\end{titlepage}
\hypersetup{pageanchor=true}
\tableofcontents

\mainmatter

% The optional short title keeps the forced line break out of the contents,
% running header, and PDF bookmarks. It can be omitted for automatic wrapping.
\chapter[Systems of Linear Equations and Matrices]{Systems of Linear\\Equations and Matrices}
\label{ch:systems}
% Place this before the first section. The list is generated automatically.
\chaptercontents

\begin{chapterintroduction}
Many problems become easier to study when their quantities are arranged in rows
and columns. A matrix keeps those quantities in order, so that the same rules
can be applied to measurements, equations, and other kinds of data. In this
chapter we use that idea to solve systems of linear equations and to explain
what their solutions represent. For example, suppose two unknowns satisfy the
equations
\[
  \begin{aligned}
    5x+y &= 3,\\
    2x-y &= 4.
  \end{aligned}
\]
The same information is stored in the augmented matrix
\(\left[\begin{smallmatrix}5&1&3\\2&-1&4\end{smallmatrix}\right]\).
\end{chapterintroduction}

% In this example the section text starts on the following page.
\clearpage

\section{Introduction to Systems of Linear Equations}
\label{sec:systems}
A system of equations collects several conditions on the same unknowns. A
solution must satisfy every equation at once. In two variables, each equation
with a nonzero coefficient describes a line. The lines may intersect once,
coincide, or be parallel and distinct. These cases correspond to one solution,
infinitely many solutions, or no solution.

\section{Gaussian Elimination}
Elimination replaces the equations with simpler ones that have the same
solutions. Adding the sample equations from the introduction gives \(7x=7\),
so \(x=1\). Substituting into either original equation gives \(y=-2\). The
same steps can be recorded as operations on the rows of a matrix.

\section{Matrices and Matrix Operations}
Matrices can be added entry by entry when they have the same dimensions.
Multiplication uses a different rule. An entry of the product comes from
pairing a row of the left matrix with a column of the right matrix and adding
the resulting products.

\section{Inverses and Algebraic Properties of Matrices}
The identity matrix leaves a matrix unchanged under multiplication wherever
the product is defined. If a square matrix \(A\) has an inverse, multiplying
by \(A^{-1}\) reverses the effect of multiplying by \(A\). For example, the
equation \(A\mathbf{x}=\mathbf{b}\) has solution
\(\mathbf{x}=A^{-1}\mathbf{b}\) when \(A\) is invertible.

\section{Elementary Matrices and a Method for Finding
  \texorpdfstring{\textcolor{ChapterAccent}{\(A^{-1}\)}}{A inverse}}
An elementary matrix is obtained by applying one elementary row operation to
an identity matrix. To find an inverse, apply the same row operations to both
blocks of \([A\mid I]\). If the left block reduces to \(I\), the right
block is the inverse.

\section{More on Linear Systems and Invertible Matrices}
An invertible coefficient matrix gives exactly one solution for every
right-hand side. When that property fails, comparing the pivots of the
coefficient and augmented matrices can help explain whether there are free
variables or incompatible equations.

\section{Diagonal, Triangular, and Symmetric Matrices}
The position of zero entries often makes a computation easier. An upper
triangular system can be solved from the last equation upward. A symmetric
matrix has matching entries across its main diagonal and frequently arises
when relationships are measured in both directions.

\section{Applications of Linear Systems}
The examples below use the same modeling procedure. Choose the unknowns,
translate each constraint into an equation, and check the solution against
the original setting.

% Each subsection appears as a bullet in the chapter contents.
\subsection{Network Analysis (Traffic Flow)}
At an intersection with steady traffic, the total flow entering equals the
total flow leaving. Writing this balance at every intersection produces a
system for the unknown road flows.

\subsection{Electrical Circuits}
An ideal circuit can be described by recording current conservation at
junctions and voltage changes around loops. The resulting equations can be
solved for the branch currents.

\subsection{Balancing Chemical Equations}
The quantity of each element is the same on both sides of a balanced
reaction. Assigning unknown coefficients to the formulas turns each element
count into a linear equation.

\subsection{Polynomial Interpolation}
Substituting data points into a polynomial with unknown coefficients gives a
system of equations. Solving it determines a polynomial that passes through
the chosen points.

\section{Leontief Input-Output Models}
When industries buy from one another, their total output must cover both
internal purchases and outside demand. If \(C\) records the input required
per unit of output, total production \(\mathbf{x}\) satisfies
\((I-C)\mathbf{x}=\mathbf{d}\), where \(\mathbf{d}\) is outside demand.
Table~\ref{tab:inputs} gives one possible matrix \(C\). Each column represents
an industry producing one unit of output.

\begin{table}[htbp]
  \centering
  \caption{Input required per unit of output}
  \label{tab:inputs}
  \begin{tabular}{@{}lrrr@{}}
    \toprule
    Input from & Agriculture & Manufacturing & Transport \\
    \midrule
    Agriculture   & 0.20 & 0.10 & 0.05 \\
    Manufacturing & 0.15 & 0.25 & 0.15 \\
    Transport     & 0.05 & 0.10 & 0.20 \\
    \bottomrule
  \end{tabular}
\end{table}

\chapter{Vectors and Linear Transformations}
\chaptercontents

% Both chapter commands accept an optional replacement heading.
\begin{chapterintroduction}[Introduction]
Vectors describe quantities such as displacement and velocity through their
components. Using the matrices from Chapter~\ref{ch:systems}, we can represent
transformations that rotate, stretch, or project vectors.
\end{chapterintroduction}

\section{Vectors in the Plane}
A vector in the plane can be represented by an ordered pair. Adding two
vectors combines the horizontal components and the vertical components
separately. Geometrically, these additions describe the total displacement
after following both vectors in sequence.

\section{Linear Combinations and Span}
A linear combination is obtained by scaling vectors and adding the results.
The span of a collection consists of all the combinations that can be
formed. Two nonparallel vectors in the plane span the entire plane.

\section{Matrix Representations}
The columns of a transformation matrix are the images of the standard basis
vectors. Once those images are known, linearity determines the image of
every vector. For instance, projection onto the horizontal axis maps
\((x,y)\) to \((x,0)\).

\end{document}
